\subsection{Solving Polynomial Inequalities in Factored Form}
The method of solving a polynomial inequality is based on analyzing the signs.
\begin{Example}
Solve the polynomial inequality
\[
\left(x-4\right)\left(2x+3\right)\left(x+5\right)\left(3-x\right)<0
\]
First, we set each factor equal to zero to find the cut-off points.
\begin{lpic}{}{2}
\end{lpic}
Mark these points (in order, but spacing doesn't matter) on a number line, and analyze each factor's sign. This is commonly done with a diagram like this one.
\begin{cpic}{}{8}
\foreach \x/\lbl in {-3/{$x-4$},-4/{$2x+3$},-5/{$x+5$},-6/{$3-x$},-8/{Product}}
	\regtextright{-7.5}{\x}{\lbl};
\draw[style=axis] (-7.5,-1) -- (7.5,-1);
\foreach \x in {-4.5,-1.5,1.5,4.5} \draw[style=tick] (\x,-.75) -- (\x,-1.25);
\end{cpic}
We are looking for the set where the polynomial is \makeblank 0, so where the product is \makeblank. So the solution set is
\[
\left\{x\suchthat\htstrut\makeblank[1in]\mbox{ or }\makeblank[1in]\mbox{ or }\makeblank[1in]\right\}
\]
or in interval notation
\[
\makeblank[1in]\cup\makeblank[1in]\cup\makeblank[1in]
\]
\end{Example}